\chapter{Allocazione del budget pubblicitario}
\label{cap:budget}

\noindent\textbf{Classe:} NLP convesso \hfill
\textbf{Script:} \texttt{python/lab08\_budget.py}

\section{Motivazione gestionale}

Come ripartire una campagna da 100.000~\euro{} tra social, search, TV e influencer,
quando ogni canale mostra rendimenti marginali decrescenti? I primi euro su un canale
rendono molto, i successivi sempre meno: la risposta è una funzione concava della
spesa, e l'ottimo segue una regola economica elegante che questo capitolo dimostra e
verifica numericamente.



L'allocazione di un budget tra impieghi con rendimenti decrescenti è uno schema
universale: marketing su più canali, fondi tra progetti, tempo tra attività,
fertilizzante tra campi. La teoria dice qualcosa di forte e verificabile:
all'ottimo i rendimenti marginali si eguagliano --- ed è esattamente ciò che il
solver restituirà, cifra per cifra.

\textbf{In questo capitolo impareremo a:} (1) modellare rendimenti decrescenti con
funzioni concave; (2) risolvere NLP convessi con scipy; (3) verificare
numericamente le condizioni KKT; (4) costruire la curva valore-budget e usare
$\lambda$ per negoziare il budget.

\section{Costruzione guidata del modello}

\paragraph{Il problema a parole.} \emph{Decidiamo} quanto spendere in ciascun canale
pubblicitario. \emph{L'obiettivo} è massimizzare la risposta totale attesa della
campagna. \emph{I vincoli}: la spesa complessiva non può superare il budget, e
ciascun canale ha un investimento massimo utile.

\subsection*{Insiemi e dati (input del modello)}
\begin{center}\small
\begin{tabular}{llp{7.2cm}}
\toprule
Simbolo & Tipo & Significato \\
\midrule
$I$ & insieme finito & canali pubblicitari \\
$b$ & $\in \R_{> 0}$ & budget totale \\
$r_i(\cdot)$ & funzione concava crescente & risposta attesa del canale $i \in I$ \\
$u_i$ & $\in \R_{> 0}$ & massimo investimento utile o consentito nel canale $i$ \\
\bottomrule
\end{tabular}
\end{center}

\subsection*{Variabili decisionali}
\begin{center}\small
\begin{tabular}{llp{7.2cm}}
\toprule
$x_i$ & $\ge 0$, continua & investimento nel canale $i \in I$ \\
\bottomrule
\end{tabular}
\end{center}

\begin{modello}[Allocazione con risposta concava]
\begin{align}
\max\;& \sum_{i \in I} r_i(x_i) \notag\\
\text{s.t.}\;\; & \sum_{i \in I} x_i \le b, \qquad 0 \le x_i \le u_i \;\;\forall i \in I \notag
\end{align}
Nel caso di studio $r_i(x) = a_i \log(1 + k_i x)$, con $a_i, k_i \in \R_{>0}$ dati:
concava, nulla in zero, con rendimento marginale
$r_i'(x) = a_i k_i / (1 + k_i x)$ decrescente.
\end{modello}

\begin{spiegazione}
\textbf{Funzione obiettivo.} La risposta totale $\sum_{i \in I} r_i(x_i)$ è separabile
per canale e concava: sommare funzioni concave preserva la concavità, quindi il
problema ha ottimo globale certificabile.

\textbf{Vincolo di budget ($\sum_{i \in I} x_i \le b$).} È l'unico vincolo che lega i
canali tra loro: senza di esso ogni canale andrebbe al proprio tetto.

\textbf{Bound ($0 \le x_i \le u_i$).} Il tetto $u_i$ modella il limite oltre cui il
canale non assorbe altra spesa utile (inventario pubblicitario, saturazione del
pubblico).
\end{spiegazione}

\begin{spiegazione}[La condizione economica all'ottimo]
Per le KKT, all'ottimo esiste $\lambda \ge 0$ (prezzo ombra del budget) tale che, per
ogni canale non bloccato ai limiti $0$ o $u_i$:
\[
r_i'(x_i^*) = \lambda .
\]
In parole: \textbf{l'ultimo euro investito rende lo stesso in tutti i canali attivi}.
Se così non fosse, spostare un euro dal canale con marginale basso a quello con
marginale alto migliorerebbe la risposta: non saremmo all'ottimo. I canali al tetto
$u_i$ possono avere marginale $> \lambda$ (vorremmo investirci di più ma non si può);
quelli a zero hanno marginale iniziale $< \lambda$ (non valgono nemmeno il primo euro).
\end{spiegazione}

\section{Esempio svolto a mano}

\begin{esempio}[Due canali logaritmici]
$r_1(x) = 100\log(1 + 0{,}2x)$, $r_2(x) = 60\log(1 + 0{,}1x)$, budget $b = 50$
(migliaia di \euro), senza tetti.

\textbf{Passo 1 — condizione di uguaglianza dei marginali.}
\[
\frac{100 \cdot 0{,}2}{1 + 0{,}2x_1} = \frac{60 \cdot 0{,}1}{1 + 0{,}1x_2}
\quad\Longrightarrow\quad
\frac{20}{1 + 0{,}2x_1} = \frac{6}{1 + 0{,}1x_2}.
\]

\textbf{Passo 2 — budget.} $x_2 = 50 - x_1$. Sostituendo:
$20\,(1 + 0{,}1(50 - x_1)) = 6\,(1 + 0{,}2 x_1)$, cioè
$20\,(6 - 0{,}1x_1) = 6 + 1{,}2x_1 \Rightarrow 120 - 2x_1 = 6 + 1{,}2x_1
\Rightarrow x_1 = \dfrac{114}{3{,}2} = 35{,}6$, $x_2 = 14{,}4$.

\textbf{Passo 3 — prezzo ombra.} $\lambda = 20/(1 + 0{,}2 \cdot 35{,}6) = 2{,}46$:
un euro di budget in più renderebbe $\approx 2{,}46$ unità di risposta ---
lo stesso numero da qualunque canale lo si faccia entrare.
\end{esempio}

\section{Caso di studio e implementazione}

Quattro canali, $r_i(x) = a_i\log(1 + k_i x)$, budget $b = 100$ (migliaia di \euro), dati in
\texttt{dati/budget\_canali.csv}. Il problema è concavo e lo risolviamo con
\texttt{scipy} con il metodo SLSQP (\emph{Sequential Least Squares Programming},
un algoritmo per NLP vincolati che risolve una successione di sottoproblemi
quadratici):

\begin{lstlisting}[style=python]
res = minimize(lambda x: -risposta(x),
               x0=np.full(4, B/4),
               bounds=[(0, ui) for ui in u],
               constraints=[{"type": "ineq",
                             "fun": lambda x: B - x.sum()}],
               method="SLSQP")
\end{lstlisting}

\section{Risultati e verifica delle KKT}

\begin{lstlisting}[style=output]
Risposta totale: 1.449,8 (migliaia di contatti utili)
     canale |  spesa | tetto | risposta | marginale
     social |   24,3 |    60 |    385,3 |   8,2693
     search |   34,8 |    80 |    539,5 |   8,2693
         TV |   22,9 |   120 |    235,2 |   8,2693
 influencer |   18,0 |    35 |    289,8 |   8,2693

Verifica: +1000 EUR di budget -> risposta +8,244 ~ lambda
\end{lstlisting}

\begin{insight}
La condizione teorica si vede nei numeri: i quattro marginali coincidono alla quarta
cifra ($\lambda = 8{,}2693$), e l'aumento reale di risposta con mille euro in più
(8,244) conferma la lettura del moltiplicatore. Nota manageriale: la TV riceve
\emph{meno} dei social nonostante il tetto più alto --- non conta la dimensione del
canale ma la velocità con cui satura ($k_i$).
\end{insight}

\begin{center}
\begin{tikzpicture}
\begin{axis}[labmezza, title={Curve di risposta concave},
  xlabel={spesa nel canale (migliaia \euro)},
  ylabel={risposta (migliaia contatti)}]
\addplot [teal, thick] table [col sep=comma, x=x, y=social] {\figdat{cap08_risposte}};
\addplot [rossomattone, thick] table [col sep=comma, x=x, y=search] {\figdat{cap08_risposte}};
\addplot [verde, thick] table [col sep=comma, x=x, y=TV] {\figdat{cap08_risposte}};
\addplot [arancio, thick] table [col sep=comma, x=x, y=influencer] {\figdat{cap08_risposte}};
\legend{social, search, TV, influencer}
\end{axis}
\end{tikzpicture}%
\begin{tikzpicture}
\begin{axis}[labmezza, title={Valore del budget e $\lambda$ decrescente},
  xlabel={budget totale (migliaia \euro)}, ylabel={risposta totale ottima}]
\addplot [teal, very thick] table [col sep=comma, x=budget, y=risposta]
  {\figdat{cap08_valore_budget}};
\legend{risposta ottima}
\end{axis}
\end{tikzpicture}
\captionof{figure}{A sinistra: le curve di risposta $r_i(x) = a_i\log(1 + k_i x)$
dei quattro canali --- concave, con velocità di saturazione diverse. A destra: la
risposta totale ottima al crescere del budget (curva valore), concava: ogni tranche
di budget rende meno della precedente.}
\end{center}

\begin{center}
\begin{tikzpicture}
\begin{axis}[lab, width=0.9\textwidth, stack plots=y, grid=none, area legend,
  title={Mix ottimo al crescere del budget: i canali saturano ai tetti},
  xlabel={budget totale (migliaia \euro)}, ylabel={spesa per canale},
  legend columns=4, legend style={at={(0.5,-0.24)}, anchor=north},
  enlargelimits=false, ymin=0]
\addplot [fill=teal!80, draw=none] table [col sep=comma, x=budget, y=social]
  {\figdat{cap08_mix}} \closedcycle;
\addplot [fill=rossomattone!75, draw=none] table [col sep=comma, x=budget, y=search]
  {\figdat{cap08_mix}} \closedcycle;
\addplot [fill=verde!75, draw=none] table [col sep=comma, x=budget, y=TV]
  {\figdat{cap08_mix}} \closedcycle;
\addplot [fill=arancio!80, draw=none] table [col sep=comma, x=budget, y=influencer]
  {\figdat{cap08_mix}} \closedcycle;
\legend{social, search, TV, influencer}
\end{axis}
\end{tikzpicture}
\captionof{figure}{Mix ottimo di spesa (aree impilate, migliaia di \euro) al crescere
del budget totale: i canali entrano e crescono insieme mantenendo uguali i ritorni
marginali; l'influencer satura per primo al suo tetto (35), seguito dagli altri.}
\end{center}

\section{Analisi di sensitività}

\begin{lstlisting}[style=output]
b =  20: risposta   515,9   lambda = 18,971
b =  60: risposta 1.070,4   lambda = 10,909
b = 100: risposta 1.449,8   lambda =  8,244
b = 180: risposta 1.988,2   lambda =  5,538
b = 260: risposta 2.370,2   lambda =  4,049
b = 300: risposta 2.498,1   lambda =  0,000  (tutti i canali ai tetti)
\end{lstlisting}

\begin{spiegazione}
La curva valore-budget è concava: $\lambda$ scende da 19 a 4 man mano che il budget
cresce. A $b = 300$ tutti i canali sono ai tetti ($60+80+120+35 = 295$): il budget
smette di essere la risorsa scarsa e $\lambda$ crolla a zero. Per un CFO: la
\emph{curva} di $\lambda$ è l'argomento quantitativo per decidere quanto budget
concedere al marketing --- si finanzia finché $\lambda$ supera il valore di un euro
investito altrove.
\end{spiegazione}

\section{Esercizi}

\begin{esercizio}[verifica]
Risolvere l'esempio a mano con $b = 80$: dove finisce il budget aggiuntivo? Come
cambia $\lambda$?
\end{esercizio}

\begin{esercizio}[risposta esponenziale]
Sostituire la risposta della TV con $r(x) = 520\,(1 - e^{-0{,}025x})$ e ripetere
l'ottimizzazione: cambia il mix? La funzione ha un tetto naturale: il bound $u_i$
serve ancora?
\end{esercizio}

\begin{esercizio}[canale escluso]
Aggiungere un quinto canale con marginale iniziale $a_5 k_5 = 5 < \lambda$: verificare
che l'ottimo lo lascia a zero e interpretare il suo ``costo ridotto''
$a_5 k_5 - \lambda$.
\end{esercizio}

\begin{esercizio}[copertura equa]
Il budget serve tre segmenti di clientela con coperture diverse per canale.
Massimizzare la copertura del segmento peggiore ($\max z$ con
$z \le \text{copertura}_g$): quanto costa l'equità in risposta totale?
\end{esercizio}
